% Coursework and industrial vision are evidence of preparation, not clinical MRI experience.
\renewcommand{\cvfootertitle}{Cover Letter}
\renewcommand{\cvBody}{%
    \cvletterhead
        {Stuttgart, 25 October 2026}
        {Chipiron\\AI Engineering team\\Paris, France}
        {Application: AI Engineer (Always Open)}

    Dear Chipiron team,

    % Earlier opening retained: Your approach brings together two parts of my background: medical-imaging coursework in electrical engineering during both my bachelor's and master's studies, and hands-on AI development at Fraunhofer IPA.
    Making MRI simpler and more accessible connects directly with my academic background and my work at Fraunhofer IPA. During my university studies, I took advanced courses in AI and computer graphics, including practical work rendering MRI images, alongside medical-imaging coursework in electrical engineering during both my bachelor's and master's studies. I would like to bring that foundation and my hands-on AI development experience to Chipiron.

    In my industrial vision work, I have owned the path from datasets to deployed models: creating and annotating data, training in PyTorch, evaluating results and integrating inference into working systems. Retrofitting a Berghof PLC with a Hailo accelerator meant making a model useful within real hardware constraints, then improving it through customer feedback. At Chipiron, I would bring that same attention to the whole imaging pipeline and work closely with acquisition and reconstruction colleagues to understand what an improvement must achieve in practice.

    Your emphasis on proving that models work particularly interests me. I co-authored \textit{Semantic Robustness Probing via Inpainting}, which uses diffusion-based inpainting to investigate object-detector robustness. This is computer-vision research rather than MRI reconstruction, but it provides a relevant starting point for thinking about evaluation beyond a single aggregate score. I would want to understand which structures must be preserved, how failures appear and how performance changes across acquisition conditions.

    I also enjoy removing the engineering obstacles that slow research down. I ported a quantum-physics simulation to CUDA and PyTorch, profiling it to achieve an approximately 700-fold speedup, and built a vision-language self-labeling tool that reduced manual labeling time by 75\%. Alongside maintaining institute repositories in Python, Rust and C++, I helped turn a self-hosted coding assistant into a service researchers use, on my own initiative with a colleague.

    I would bring practical model-development experience, curiosity about the imaging physics and a willingness to take ownership across disciplinary boundaries. My medical-imaging background is academic; I am keen to develop MRI-specific expertise alongside your team. I work in English at C1 level, speak French at B1 and am open to relocating to Paris. I would welcome a conversation about where my experience could be most useful as Chipiron grows.

    \cvletterclosing{Best regards}
}
